\section{The Max-Cut reduction}\label{sec:reduction}

Fix a rational $t\in(0,1)$, tree parameters $s,m$, a rational
$0<\mu<1/2$, and $r=1-\mu/d$. The parameters will be chosen in
\Cref{sec:completion}. Let $\mathcal G$ be an affine game as in
\Cref{prop:affine-games}, with edge law $\lambda$ and alphabet dimensions
$n+1,n$. Let $T\ge1$ be an integer. We use the product label spaces
\[
 \mathcal B=(\F^{n+1})^T,\qquad
 \mathcal A=(\F^n)^T,
\]
and the row space $H=\Aff(\mathcal A,\F)$.

\subsection{Vertices and edge distribution}

For each $\underline x\in X^T$, make one graph vertex
$(\underline x,w)$ for every word $w\in\bits^{\mathcal B}$.
This is the standard long-code word domain~\cite{BGS98,Hastad01};
the tree construction specifies how the test samples its queries.
Thus an arbitrary cut is exactly a family of functions
\[
 \sigma_{\underline x}:\bits^{\mathcal B}\longrightarrow\bits.
\]
The functions receive only the resulting words; they do not receive the
rows, gate tables, slice, or sampled game edges that produced them.

The graph's edge distribution is the following test.
\begin{enumerate}
\item Draw $\underline y\in Y^T$ from the product right-endpoint marginal
of $\lambda$. Independently draw two edge tuples
$\underline e,\underline e'$ from the product conditional edge law given
$\underline y$. Write $\underline x,\underline x'$ for their left
endpoints and
$\pi_{\underline e},\pi_{\underline e'}:\mathcal B\to\mathcal A$
for their coordinatewise projections.

\item Independently of all game data, draw a random slice $\mathcal D$.
Draw $F,F'$ as in \Cref{lem:tree-correlation}: corresponding entries
are independent $t$-correlated uniform pairs on $\mathcal D$, and
identical uniform pairs off $\mathcal D$.

\item Draw a uniform common row array $R\in H^k$. In each of two copies,
independently retain each row with probability $r$ and otherwise resample
it uniformly, obtaining $R^1,R^2$.

\item Output the unordered edge whose endpoints are
\begin{equation}\label{eq:test-queries}
 \bigl(\underline x,W(R^1,F)\circ\pi_{\underline e}\bigr),
 \qquad
 \bigl(\underline x',-W(R^2,F')\circ\pi_{\underline e'}\bigr).
\end{equation}
\end{enumerate}
All edge weights are the probabilities of these outcomes, aggregated over
outcomes with the same unordered endpoints.

\begin{lemma}[Completeness]\label{lem:completeness}
If $\Val(\mathcal G)\ge1-a$, the graph has a cut of value at least
\begin{equation}\label{eq:completeness}
 (1-2Ta)\frac{1+t r^{2d}}2
 \ge\frac{1+t}{2}-t\mu-2Ta.
\end{equation}
The first lower bound is useful when $2Ta\le1$, which is the regime used
below.
\end{lemma}

\begin{proof}
Choose game labels $\ell_X,\ell_Y$ satisfying an edge with probability at
least $1-a$. Label each word vertex by evaluation at the tuple of left
labels:
\[
 \sigma_{\underline x}(w)=
 w\bigl(\ell_X(x_1),\ldots,\ell_X(x_T)\bigr).
\]
Each coordinate of either sampled edge tuple has the original edge law.
A union bound shows that all $2T$ sampled constraints are satisfied except
with probability at most $2Ta$. When all these constraints are satisfied,
the two queried tuples
project to the same label
$l=(\ell_Y(y_1),\ldots,\ell_Y(y_T))\in\mathcal A$.
The cut signs then have product
$-W(R^1,F)(l)W(R^2,F')(l)$. The word-generation data are independent of
the game edges, so \Cref{lem:tree-correlation} gives conditional cut
probability $(1+t r^{2d})/2$. On the remaining event the cut probability
is nonnegative. This proves the first bound when $2Ta\le1$; if $2Ta>1$
that bound is negative and is automatic. Finally,
$r^{2d}\ge1-2\mu$, and the loss due to the prefactor is at most $2Ta$,
giving the second bound.
\end{proof}

\subsection{A large cut produces excess stability}

An arbitrary cut need not change sign when its queried word is negated.
We separate its two parities because the negative second query makes
odd dependence increase the cut value and even dependence decrease it.
For each cut table define its even and odd parts by
\[
 \sigma^{\mathrm{even}}(w)=\frac{\sigma(w)+\sigma(-w)}2,
 \qquad
 \sigma^o(w)=\frac{\sigma(w)-\sigma(-w)}2.
\]
For an edge tuple and a right-endpoint tuple, set
\begin{align}
 G_{\underline e}(R,F)
 &=\sigma^o_{\underline x(\underline e)}
       (W(R,F)\circ\pi_{\underline e}),\label{eq:edge-function}\\
 G_{\underline y}(R,F)
 &=\E_{\underline e\mid\underline y}G_{\underline e}(R,F),
 \qquad \bar G_{\underline y}=T_rG_{\underline y}.
 \label{eq:average-function}
\end{align}
Write $A_{\underline y}$ and $\bar A_{\underline y}=T_rA_{\underline y}$
for the analogous averages using $\sigma^{\mathrm{even}}$.
These functions are bounded in absolute value by one.
Moreover,
\[
 G_{\underline y}(R,F)=\phi_{\underline y}(W(R,F)),\qquad
 \phi_{\underline y}(w)=\E_{\underline e\mid\underline y}
 \sigma^o_{\underline x(\underline e)}(w\circ\pi_{\underline e}),
\]
where $\phi_{\underline y}$ is odd and depends on no slice or row data.
Thus $\bar G_{\underline y}$ has mean zero in the slice entries for
every fixed $\underline y,\mathcal D,R,F_{\mathrm{out}}$, by the
observation following \Cref{lem:tree-correlation}.
The later extraction argument, \Cref{prop:extraction}, applies to precisely
such bounded odd functions of the generated word.

\begin{lemma}[Stability identity]\label{lem:cut-stability}
For every cut of the constructed graph,
\begin{equation}\label{eq:cut-stability}
 \cut(\sigma)=\frac12+
 \frac12\E\Stab_t^{F_{\mathcal D}}(\bar G_{\underline y})
 -\frac12\E\Stab_t^{F_{\mathcal D}}(\bar A_{\underline y}),
\end{equation}
where both expectations are over
$\underline y,\mathcal D,R,F_{\mathrm{out}}$.
In particular,
\begin{equation}\label{eq:cut-upper-stability}
 \cut(\sigma)\le\frac12+
 \frac12\E\Stab_t^{F_{\mathcal D}}(\bar G_{\underline y}).
\end{equation}
\end{lemma}

\begin{proof}
Condition on $\underline y,\mathcal D,R,F_{\mathrm{out}}$.
The sampled edge tuples and the two row resamplings are independent
between copies. After averaging over them, the expected first sign as a
function of $F_{\mathcal D}$ is $\bar A_{\underline y}+\bar G_{\underline y}$.
The minus sign in the second queried word makes the expected second sign
$\bar A_{\underline y}-\bar G_{\underline y}$ at $F'_{\mathcal D}$.
The two slice inputs are ordinary $t$-correlated uniform cube inputs.
The rows in these functions are the common latent array $R$; the two
independent row resamplings have already been absorbed into $T_r$.
Negating every slice entry leaves $\bar A_{\underline y}$ unchanged and
negates $\bar G_{\underline y}$. Their Fourier supports therefore have
opposite total parities, so the cross terms vanish. The expected sign
product is
$\Stab_t(\bar A_{\underline y})-\Stab_t(\bar G_{\underline y})$.
An edge is cut with probability one half of one minus its sign product,
giving \eqref{eq:cut-stability}. Nonnegativity of stability for $t>0$
gives \eqref{eq:cut-upper-stability}.
\end{proof}

If a cut has value at least $(1+B(t))/2+\delta/2$, then its odd average
has expected slice stability at least $B(t)+\delta$.
The next three sections show that this excess stability yields a label
carrying substantial Fourier mass, with constants independent of the game
dimension and the number of repetitions; the formal statement is
\Cref{prop:extraction}. \Cref{sec:decoding} then turns that mass into a
labeling of a single copy of the original game.
